% LaTeX source maintained directly for Overleaf.
\chapter{VI. Ứng dụng và vận hành hệ thống}

\section{Stakeholder và use case}

\begin{itemize}
\item Người dùng cuối hỏi thông tin phim và nhận recommendation có lý do.
\item Data engineer thu thập, chuẩn hóa, enrich và import dữ liệu.
\item Người vận hành khởi động Neo4j/API và kiểm tra health/stats.
\item Nhà phát triển và nghiên cứu viên chạy test, query, benchmark và truy vết evidence.
\end{itemize}

\section{Yêu cầu chức năng}

\begin{longtable}{|>{\raggedright\arraybackslash}p{0.29\textwidth}|>{\raggedright\arraybackslash}p{0.29\textwidth}|>{\raggedright\arraybackslash}p{0.29\textwidth}|}
\caption{Tổng hợp yêu cầu chức năng}\label{tab:6-1} \\
\hline
\textbf{ID} & \textbf{Yêu cầu} & \textbf{Artifact chính} \\ \hline
\endfirsthead\hline \textbf{ID} & \textbf{Yêu cầu} & \textbf{Artifact chính} \\ \hline\endhead
F01 & Thu thập/cache TMDB & \texttt{collect\_tmdb.py}, \texttt{tmdb\_client.py} \\ \hline
F02 & Tải/stream IMDb ratings & \texttt{download\_imdb.py}, \texttt{imdb\_loader.py} \\ \hline
F03 & Clean, normalize, entity resolution & \texttt{processing/} \\ \hline
F04 & Sinh 5 node CSV + 5 edge CSV + manifest & \texttt{pipeline.py} \\ \hline
F05 & Import/validate Neo4j idempotent & \texttt{load\_neo4j.py} \\ \hline
F06 & RDF export, entailment, 10 SPARQL & \texttt{export\_rdf.py}, \texttt{semantic\_reasoning.py} \\ \hline
F07 & 10 Cypher và CRUD parameterized & \texttt{cypher/}, \texttt{kg/crud.py} \\ \hline
F08 & QA answer + evidence & \texttt{/ask} \\ \hline
F09 & Recommendation + explanation & \texttt{/recommend} \\ \hline
F10 & Search/stats/health và UI & FastAPI/static UI \\ \hline
\end{longtable}

\section{Hỏi–đáp}
\reportfigure{qa_sequence.pdf}{Luồng xử lý một yêu cầu hỏi--đáp}{fig:qa-sequence}

\begin{Verbatim}[fontsize=\scriptsize,frame=single]
Question → parser 9 intent → trích xuất slot → entity linking
         → fixed query catalog → parameterized Cypher
         → Neo4j → answer + evidence + latency
\end{Verbatim}

Parser tất định nhận diện chín intent và trích xuất các slot thực thể hoặc ngưỡng
rating. Entity linker canonicalize tên trước query và trả confidence trong
evidence. Mỗi intent ánh xạ tới một graph pattern cố định trong query catalog;
user input luôn nằm trong parameter. Câu hỏi ngoài phạm vi trả intent
\texttt{unknown} thay vì tạo truy vấn tùy ý.

\subsection{Giá trị của lớp hỏi--đáp đối với Cypher}

Lớp hỏi--đáp không thay thế Cypher mà cung cấp một giao diện có kiểm soát cho
query catalog. Mỗi intent chọn một graph pattern cụ thể; entity linker bổ sung
stable ID; repository truyền slot qua parameter và giữ nguyên cấu trúc truy vấn.
Nhờ đó, người dùng không cần viết Cypher nhưng kết quả vẫn thể hiện trực tiếp các
năng lực traversal, aggregation, variable-length path và truy vấn fact suy ra
của Neo4j.

\begin{longtable}{|>{\raggedright\arraybackslash}p{0.22\textwidth}|>{\raggedright\arraybackslash}p{0.28\textwidth}|>{\raggedright\arraybackslash}p{0.20\textwidth}|>{\raggedright\arraybackslash}p{0.20\textwidth}|}
\caption{Ánh xạ intent hỏi--đáp sang năng lực Cypher}\label{tab:6-2} \\
\hline
\textbf{Intent} & \textbf{Graph pattern chính} & \textbf{Năng lực Cypher} & \textbf{Evidence} \\ \hline
\endfirsthead\hline \textbf{Intent} & \textbf{Graph pattern chính} & \textbf{Năng lực Cypher} & \textbf{Evidence} \\ \hline\endhead
\texttt{movies\_by\_director} & Person-\texttt{DIRECTED}-Movie & pattern matching, parameter & Movie ID/title và relationship \\ \hline
\texttt{movies\_by\_person} & Person-\texttt{ACTED\_IN|DIRECTED}-Movie & relationship alternative, \texttt{DISTINCT} & loại vai trò thực tế \\ \hline
\texttt{actors\_in\_movie} & Person-\texttt{ACTED\_IN}-Movie & edge-property projection, sort & character và cast order \\ \hline
\texttt{common\_movies} & Person-Movie-Person & multi-hop shared neighbor & phim chung \\ \hline
\texttt{movies\_by\_genre\_rating} & Movie-\texttt{HAS\_GENRE}-Genre & filter, \texttt{coalesce}, ordering & rating được dùng \\ \hline
\texttt{co\_stars} & Person-\texttt{CO\_STARRED\_WITH}-Person & derived-edge traversal & movie count và evidence IDs \\ \hline
\texttt{directors\_by\_genre} & Person-Movie-Genre & multi-hop aggregation, \texttt{count/collect} & số lượng và Movie IDs \\ \hline
\texttt{shortest\_path} & Person-\texttt{[*..8]}-Person & variable-length shortest path & node và relationship path \\ \hline
\texttt{similar\_movies} & Movie-feature-Movie & candidate traversal và ranking & shared feature và contribution \\ \hline
\end{longtable}

Ví dụ, câu hỏi “Đường liên hệ giữa Leonardo DiCaprio và Christian Bale?” không
được trả lời bằng luật chuỗi trong Python. Parser chỉ chọn
\texttt{shortest\_path}; hai tên được liên kết về \texttt{person\_id}; sau đó
Neo4j thực thi \texttt{shortestPath((a:Person)-[*..8]-(b:Person))}. Response trả
lại chính danh sách node và relationship trên đường tìm được. Vì vậy impact của
tính năng nằm ở việc đưa các truy vấn Cypher nhiều bước đến người dùng cuối mà
không mất parameterization và khả năng truy vết evidence.

\section{Recommendation}
\reportfigure{recommendation_explanation.pdf}{Luồng xếp hạng và tạo giải thích gợi ý phim}{fig:recommendation-flow}

Neo4j tạo candidate bằng các phim chia sẻ ít nhất một director, actor, keyword,
genre hoặc studio. Mỗi feature được tính document frequency trong graph và đóng
góp IDF theo type weight. Query trả candidate, score và feature chung; Python chỉ
định dạng \texttt{Recommendation} và explanation. Tie-break theo title giúp kết quả xác
định. Endpoint nhận \texttt{movie\_id} thay vì title để tránh nhập nhằng.

\section{API và UI}
\reportfigure{web_ui.pdf}{Giao diện hỏi--đáp và gợi ý phim}{fig:web-ui}

\begin{longtable}{|>{\raggedright\arraybackslash}p{0.44\textwidth}|>{\raggedright\arraybackslash}p{0.44\textwidth}|}
\caption{Tổng hợp api và ui}\label{tab:6-3} \\
\hline
\textbf{Endpoint} & \textbf{Chức năng} \\ \hline
\endfirsthead\hline \textbf{Endpoint} & \textbf{Chức năng} \\ \hline\endhead
GET \texttt{/health} & kiểm tra Neo4j connectivity \\ \hline
GET \texttt{/stats} & số node theo label và tổng relationship \\ \hline
GET \texttt{/entities/search} & full-text rồi fallback parameterized search \\ \hline
POST \texttt{/ask} & QA answer, intent, evidence, query time \\ \hline
POST \texttt{/recommend} & top-K recommendation và explanation \\ \hline
\end{longtable}

UI có hai tab QA và recommendation, autocomplete phim, lịch sử hiển thị và state
loading/error/success. Lịch sử chỉ ở frontend; backend hiện stateless.

\section{An toàn và vận hành}

Secret nằm trong \texttt{.env}; raw data không commit. Public API không có write
endpoint. Hạn chế hiện tại
là chưa có authentication/rate limiting, nên cấu hình phù hợp demo local chứ
không phải public production deployment.

\section{Failure modes}

\begin{longtable}{|>{\raggedright\arraybackslash}p{0.44\textwidth}|>{\raggedright\arraybackslash}p{0.44\textwidth}|}
\caption{Tổng hợp failure modes}\label{tab:6-4} \\
\hline
\textbf{Tình huống} & \textbf{Hành vi} \\ \hline
\endfirsthead\hline \textbf{Tình huống} & \textbf{Hành vi} \\ \hline\endhead
Câu hỏi ngoài 9 intent & trả intent \texttt{unknown} và gợi ý dạng câu hỏi hỗ trợ \\ \hline
Thiếu slot bắt buộc & không chạy Cypher tùy ý; trả không hiểu/không tìm thấy \\ \hline
entity không tồn tại & trả entity-not-found/không tìm thấy \\ \hline
movie ID recommend sai & HTTP 404 \\ \hline
Neo4j không sẵn sàng & health 503 \\ \hline
tên fuzzy nhập nhằng & không tự chọn tied candidate \\ \hline
\end{longtable}

Failure behavior là một phần correctness. Hệ QA an toàn nên thừa nhận không biết
hoặc yêu cầu làm rõ thay vì chạy query với entity đoán sai.

\section{Tính faithful của explanation}

QA evidence gồm canonical entity/link confidence và record/path trả từ Neo4j.
Recommendation explanation không được model sinh: nó tổng hợp từ shared feature
đã đóng góp score. Cách làm giảm độ tự nhiên so với NLG explanation \cite{colas2023knowledge} nhưng
có tính faithful cao trong phạm vi hệ thống: mỗi lý do tương ứng quan hệ tồn tại
trong graph và contribution xác định.
